\section{Implementation in \fevc{} 0.6.0}\label{sec:impl}

The prototype behind Sections~\ref{sec:q1} and~\ref{sec:evidence} preceded the package. \fevc{} implemented the
method between November 13 and 14, 2026, and released it in 0.6.0 as its only standard-error method and its
default, \code{inference(auto)}; \code{inference(none)} gives point estimates only. What the command does is
specified in the package's \code{fevc/docs/INFERENCE.md}; where this note and that guide differ, the guide
describes the command. This section lists the differences, states the rule by which the default chooses an
interval, and summarises the package's own Monte Carlo evidence.

\subsection{Differences from the prototype}\label{sec:impl-diff}

\begin{enumerate}[leftmargin=1.6em,itemsep=3pt]
\item \textbf{Deletion modes and repeat-measurement sets.} The propositions and the prototype cover match
deletion, where stayers are the only sets of rows with identical regressors (Section~\ref{sec:scope}); the control
term of Section~\ref{sec:controls} is derived for both modes. 0.6.0 also covers observation deletion,
where every worker--firm cell with two or more observations, a mover's match included, is a set, with the pair
weight $B_{ii}/(1-P_{ii})$ for every target; and mover-only populations. The only sets are deletion units that
share a worker--firm cell. Co-mover groups and cycles of matches, which earlier drafts considered, are not used.
\item \textbf{Fitting the model for the other units.} The two steps and the floors are those of
Section~\ref{sec:transfer}. In addition, each step keeps at most one coefficient per hundred fitted units,
shortening the basis in a fixed order; a term that the others span is dropped, also in a fixed order, where
Section~\ref{sec:transfer} withholds the standard error; with fewer than a hundred units outside sets the model is fitted on all units
of movers; the final floor also applies to the sample variance of a set; and the pass stops with
\code{JLA\_CONSTRAINT\_FAILED} when the mean leave-out estimate is not positive (as in Section~\ref{sec:transfer}) or when every
fitted variance is floored.
\item \textbf{Leverages, projections and draws.} The prototype uses the exact $P_{ii}=1/T_g$ and
$B_{ii}=A_{gg}/T_g^2$ for stayers and 1{,}600 projections for movers, without a finite-projection correction
(Section~\ref{sec:computation}). 0.6.0 takes them from the route: on the native route the pass uses the JLA
leverages of the point estimate (200 projections by default) and checks them against its bias correction; the
Mata \code{algorithm(jla)} stage draws its own projections; exact leverages exist only on the Mata
\code{algorithm(exact)} route. The stayers' exact values are not substituted. The
default is eight pseudo-noise draws (\code{inferencedraws()}), not 16, and the linear systems are solved
iteratively, not by a factorisation.
\item \textbf{Kurtosis.} The pair term is the one for errors with the kurtosis of the normal distribution; the
variant for another stated kurtosis (Appendix~\ref{app:proofs}) is not offered.
\item \textbf{No projection of the covariance.} The prototype's simulation of Section~\ref{sec:sim} set negative
eigenvalues of the estimated covariance of the three base components to zero. 0.6.0 reports the covariance as
computed; a variance estimate that is not positive withholds the normal interval of that component, and the
default then reports its $q=1$ interval where it exists.
\item \textbf{Controls.} The prototype treats the control coefficients as known. 0.6.0 includes their estimation
to first order, as derived in Section~\ref{sec:controls}: the variance of the fixed-offset estimate under the noise
of the control-adjusted outcome, on every route. Not covered: the joint estimate's additional term in the control
index (Section~\ref{ctl:joint}), error variances that move with a control inside a cell (item~7 of
Section~\ref{ctl:est}), the within-unit covariance that the pseudo-noise omits (item~1), and poorly identified
coefficients. The repair $\hat Q_{\rm full}$ of item~1 is not implemented.
\item \textbf{Centring.} The propositions and the prototype treat the sample mean used for centring as a constant. Under
\code{centering(mean)} the package centres the pseudo-noise as the outcome, so no fixed-mean approximation is made;
under \code{centering(corrected)} the standard error is that of the mean-centred estimate, with the increment
held fixed (Section~\ref{sec:centring}).
\item \textbf{Leading direction.} The prototype uses power iteration (stopped at 200 steps; one Veneto case did
not converge). 0.6.0 uses a thick-restart Lanczos iteration in the $S_{xx}$ inner product (basis six, three Ritz
vectors kept, stopped at a relative residual of $10^{-3}$ or after 40 solves, each search started from the
previous target's direction), and a dense eigendecomposition on the exact Mata route
(Appendix~\ref{kss:app:lanczos}). The draws of the pass are reused for the remainder: about 90 solves in all with
the default eight draws, where the pass alone takes 35.
\item \textbf{One interval per component.} Section~\ref{sec:q1} gives the $q=1$ interval for designs with one weak direction
and a screening value without a validated threshold, and leaves the choice between the two intervals open.
0.6.0 chooses per component by the rule of Section~\ref{sec:impl-rule}.
\item \textbf{The other variance models.} Appendix~\ref{app:cost} compares the cost with the KSS Matlab package,
with \fevc's structured route under \code{inference(highrank)}, and with the leave-out residual standard error that
was the development default at the time (never released). 0.6.0 removed those models (\code{target\_lowess},
\code{structured\_common}, \code{structured\_leverage}, \code{leaveout\_residual}); \code{inference(highrank)} now
selects the normal intervals of this method.
\end{enumerate}

The prototype's failed tolerances stand: in the validation of Section~\ref{sec:validation} the standard error of the
firm variance in the three-year window was up to 16\% too large, beyond the registered 15\%, and the $q=1$
confirmation of Section~\ref{sec:q1sim} missed its length tolerance in one design without a bottleneck.

\subsection{The interval of every component}\label{sec:impl-rule}

Under \code{inference(auto)} a component receives the $q=1$ interval of Section~\ref{sec:q1} when its leading
direction carries more than ten per cent of the estimated sampling variance,
\[
\frac{4\lambda_1^2\max(\hat b_1^{\,2}-V_b,0)\,V_b+2\lambda_1^2V_b^2}{\hat\V[\hat\theta]},
\]
with $V_b$ the estimated variance of $\hat b_1$, or when the high-rank variance estimate \eqref{eq:Vhat} is not
positive; otherwise it receives the normal interval. The variance of the sum, whose non-zero eigenvalues are all
equal (Section~\ref{sec:q1}), keeps the normal interval. A component with the $q=1$ interval has no standard error,
and \code{e(V)} is posted only when every component has a normal interval; it is then the four-by-four covariance
of the components, mapped from the three base components. \code{inference(highrank)} asks for normal intervals
only and \code{inference(q1)} for the $q=1$ interval wherever it exists; both are strict, while the default
returns point estimates with a reason when the standard errors cannot be computed. The share is a ratio of two
estimates and can exceed one. The critical value of the $q=1$ interval is simulated with 100{,}000 draws.

The statistic and the threshold were chosen afterwards, on simulations through the command (\code{dev1} below),
confirmed on new designs (\code{conf2}), and the threshold was fixed by the owner on November 14, 2026. Over the
96 cases of \code{dev1}, the rule's lowest coverage is 0.937, against 0.940 with a threshold of 5\% and 0.935 with
20\%; it reports the $q=1$ interval for 4.3\% of the components without a bottleneck and posts \code{e(V)} in 89\%
of those samples, against 9.8\% and 76\% at 5\%. On three panels of the Veneto data (2.3 to 5.2 million
person-years, both deletion modes) every component receives the normal interval under this rule (shares 0.03\% to
8.1\%), and the default costs 24\% to 30\% more time than the point estimates alone (16 threads). Neither KSS nor
this note has a theoretical result for this data-dependent choice, and two or more dominant directions are not
covered.

\subsection{The package Monte Carlo evidence}\label{sec:impl-mc}

The package's own Monte Carlo passes the plain Stata command, with a plugin built from a named commit, through
two-block synthetic panels of about 48{,}000 person-years, from no bottleneck to a strong one, with error variances
that follow or violate the fitted model and normal or skewed errors, under match and observation deletion. Each
run's manifest was frozen before any task ran. The harness is \code{fevc/tools/default\_interval/} and the records
are in \code{fevc/docs/default\_interval/}. Table~\ref{tab:impl-mc} summarises the runs that decide the rule and
its confirmation.

\begin{table}[H]
\centering\small\setlength{\tabcolsep}{4pt}
\caption{Package Monte Carlo runs of the default interval}\label{tab:impl-mc}
\begin{tabular}{@{}lllccccc@{}}
\toprule
 & & & Designs & & & \multicolumn{2}{c}{Coverage} \\
\cmidrule(l){7-8}
Run & Kind & Commit & $\times$ 2 modes & Fits & Verdict & lowest & mean \\
\midrule
\code{dev1} & development & \code{250a760c} & 16 & 32{,}000 & none & 0.937$^a$ & -- \\
\code{conf1} & confirmation & \code{fb49e16d} & 20 & 40{,}000 & FAIL & 0.937$^b$ & 0.960$^b$ \\
\code{conf2} & confirmation & \code{fb49e16d} & 20 & 40{,}000 & PASS & 0.930 & 0.9593 \\
\code{dev3\_controls} & development & \code{02d7955f} & 13 & 13{,}000 & none & 0.918 & 0.962 \\
\code{conf3\_controls} & confirmation & \code{f3938986} & 20 & 40{,}000 & PASS & 0.923 & 0.9602 \\
\bottomrule
\multicolumn{8}{@{}p{0.92\textwidth}@{}}{\rule{0pt}{2.6ex}Coverage is that of the reported interval over the
cases (design, deletion mode and component); a
replication without an interval counts as not covered. $^a$~The 10\% rule evaluated on the draws of \code{dev1},
which ran with the first threshold of 5\%. $^b$~The verdict failed on the inventory: one fit returned point
estimates (the documented fallback after \code{JLA\_CONSTRAINT\_FAILED}), the harness of the time rejected its
task, and the verdict stands as registered. The numbers are for information only and count that fit as a
replication without an interval. \code{dev1},
\code{conf1} and \code{conf2} have no controls; \code{dev3\_controls} and \code{conf3\_controls} have controls in
every design.}
\end{tabular}
\end{table}

\code{conf2} met all of its tolerances: none of its 40{,}000 fits failed or fell back to point estimates, at most 6
of the 1{,}000 replications of a case had no interval, the lowest coverage is that of the variance of the sum under
match deletion in a design with 1{,}000 bridging workers, and without a bottleneck the reported interval is at most
1.059 times as long as the normal one and \code{e(V)} is posted in 89\% of the samples. The same replications
judged with the normal interval alone have a lowest coverage of 0.789. \code{conf3\_controls} ran the same twenty
design shapes, replications and tolerances with new seeds and controls in every design, in a fixed rotation of
four kinds; its lowest cases are the worker and firm variances (0.923 and 0.926) in a design whose blocks differ,
with the control loaded on the worker effect, where their standard errors are 8\% and 7\% short. A compatibility
review carries its claim from \code{f3938986} to the original 0.6.0 candidate build commit \code{d4b2d942}: 36
default requests give byte-identical results. The released 0.6.0 plugins were built later, from \code{2082f209};
the review does not extend the simulation claims to the subsequent Stata~15 compatibility changes or to the
revised exact eigensolver, so these confirmations are of pre-release builds, not of the released binaries. Before the control term, \code{dev2\_controls} found the standard error of the variance of
the sum 9\% to 16\% too small with one adverse control; on the same outcome draws with the term of
Section~\ref{sec:controls}, \code{dev3\_controls} gives 0.97 to 1.04 of the spread and coverage of 0.938 to 0.956.
Its lowest case (the covariance, with a control index as large as the effects and no bottleneck, under match
deletion) has a standard error 12\% too small; under \code{nuisance(fixedoffset)} the same design gives, under
match deletion, 0.99 to 1.07 of the spread and coverage of 0.948 to 0.962 (\code{dev4\_fixedoffset}), so the
shortfall belongs to the joint estimate's additional term of Section~\ref{ctl:joint}, not to the standard error of
the fixed-offset estimate. (\code{deletion(observation)} with \code{nuisance(fixedoffset)} is not covered by the
standard errors: the command returns point estimates only, as all 1{,}000 such fits of \code{dev4\_fixedoffset}
did.) The same term shows where the two blocks of firms differ by three log points and the control loads on the
firm effect: under the default the standard error of the worker variance, and in two of three designs that of the
firm variance, is 0.87 to 0.94 of the spread (\code{dev6\_shapes}, \code{dev7\_adv}), and under the fixed offset,
with match deletion, 0.92 to 0.99 (\code{dev7\_fixedoffset}). Not covered by
\code{conf3\_controls}, by design: a control index as large as the effects, error variances that move with a
control inside a cell, and a control loaded on the firm effect where the blocks differ. This is simulation
evidence for designs with one bottleneck at most and independent errors; it is not a coverage guarantee, and it
does not cover two or more dominant directions.
